\section{Test Structures and Design Search}
\label{sec:structures}

\subsection{Candidate structures}

All structures share a \qty{220}{nm} silicon film with oxide below and an oxide cladding above. Strips are fully
etched, with bottom widths of 400, 450, 500, 550 and \qty{600}{nm}, each in TE and TM, which gives ten strip
candidates. Ribs are TE only, with a slab of thickness $s$ left by a partial etch of depth $d=t-s$. At $s=70$ and
\qty{90}{nm} the candidates are rib widths of 300, 400, 500 and \qty{600}{nm}. At $s=30$, 40, 50, 110, 130 and
\qty{150}{nm} they are widths of 300 and \qty{500}{nm}, which is enough to compare slabs for a two-width set. Each
structure is read out by a ring resonator, giving $\neff$ and $\ngr$. In the sets labelled ``ring and MZI'' an MZI of
the same waveguide adds one phase per structure.

\subsection{Search}

A candidate set is a subset of the strips together with ribs from one partial etch, that is, from a single slab
thickness, since a test chip normally carries one partial-etch step. Sets contain at most eight structures. The
objective is the CRLB $\sth$ of the strip angle at stage~D with realistic priors (Section~\ref{sec:theory}), and we
minimize it by exhaustive enumeration of all admissible sets of each size. Ring-only searches also require at least
one strip and one rib. A greedy search that adds one structure at a time serves as a check on the enumeration.

\subsection{Counting unknowns}

Identifiability follows from a simple count. Strips share $w$, $t$ and $\theta$, and every strip adds two observables
but no unknowns. Their angle and width columns are parallel (Section~\ref{sec:theory}), however, so they constrain
only one combination of the two. A rib from a separate etch adds two observables ($\neff$, $\ngr$) and two unknowns
($d$, $\theta_r$). Its observables are used up determining its own etch, and nothing is left for $\theta$. A second
width of the same rib adds two observables and no unknowns. The two widths sense $d$ and $\theta_r$ with different
weights, which frees the rib's width information to break the strip degeneracy. Stage~B, where the rib shares the
strip angle, is the exception, because one rib then adds only one unknown.

\subsection{What the MZI adds}

An MZI phase is proportional to $\neff$, so its row in $\mathbf{J}$ is parallel to the ring's $\neff$ row and only
lowers the effective $\neff$ noise. For the two sets discussed in Section~\ref{sec:results}, the CRLB with ring and MZI
equals the CRLB of the ring alone with the combined $\neff$ noise to better than one part in $10^{12}$. Because the
two-width set gains about \qty{1}{\percent} from the MZI (Section~\ref{sec:mzi}), we propose ring-only test
structures. Several results below were computed with the MZI included, and we label them as such.
